\subsection{The ZTF Bright Transient Survey}
\label{sec:data:bts}

Our dataset is drawn from the Zwicky Transient Facility Bright Transient 
Survey \citep[BTS;][]{Fremling2020, Perley2020}, a magnitude-limited 
($m < 18.5$) spectroscopic completeness survey operating since 2018. 
The BTS provides spectroscopic classifications for ZTF transients, 
making it the largest homogeneous sample of spectroscopically confirmed 
extragalactic transients currently available. We use spectroscopic 
classifications as ground truth labels throughout this work.

\subsection{Sample Construction}
\label{sec:data:sample}

We queried the ALeRCE broker API \citep{Forster2021} for classification 
probabilities corresponding to all BTS objects with confirmed 
spectroscopic types on 22 February 2026. ALeRCE's light curve classifier 
\citep{Sanchez-Saez2021} outputs probabilities across 15 classes derived 
from features of the ZTF alert photometric time series. For each object 
we retrieved the most recent classification available at query time, 
corresponding to classifier versions \texttt{24.4.1}--\texttt{27.5.7a32.dev1} 
active in the production pipeline. We provide the full list of ZTF object 
identifiers and retrieved probability vectors as a machine-readable 
supplement to this paper.

We map the 11 BTS spectroscopic types to the 5 ALeRCE transient classes 
most relevant to extragalactic time-domain science: Type Ia supernovae 
(SNIa), Type II supernovae (SNII), stripped-envelope supernovae (SNIbc), 
superluminous supernovae (SLSN), and tidal disruption events (TDE). 
The full mapping is given in Table~\ref{tab:mapping}.

Of 1,987 BTS objects queried, 1,606 returned valid ALeRCE classifications 
(80.1\% success rate). Failed queries reflect objects observed before 
ALeRCE's alert ingestion began or with insufficient photometric history 
for classification. The final sample comprises 757 SNIa, 560 SNII, 
204 SNIbc, 55 SLSN, and 30 TDE.

To avoid class imbalance biasing our rarest-class estimates, we collected 
all available BTS objects in the SLSN and TDE categories rather than 
applying a fixed per-class cap. This yields a sample in which common 
classes (SNIa, SNII) are well-represented while preserving the full 
available statistics for scientifically rare but high-priority classes.

\begin{deluxetable}{ll}
\tablecaption{BTS spectroscopic type to ALeRCE class mapping.
\label{tab:mapping}}
\tablehead{\colhead{BTS Types} & \colhead{ALeRCE Class}}
\startdata
SN~Ia, SN~Ia-91T, SN~Ia-91bg, SN~Ia-pec & SNIa \\
SN~II, SN~IIn, SN~IIb, SN~IIP           & SNII \\
SN~Ib, SN~Ic, SN~Ib/c, SN~Ic-BL         & SNIbc \\
SLSN-I                                   & SLSN \\
TDE                                      & TDE \\
\enddata
\end{deluxetable}

\subsection{Fink Broker Data}
\label{sec:data:fink}

We additionally queried the Fink broker \citep{Moller2021} for the same 
set of ZTF objects via the Fink Portal API on 24 February 2026, using 
HTTP POST requests to \texttt{api.fink-portal.org/api/v1/objects}. 
Fink provides two binary classifiers relevant to our analysis. 
\texttt{snn\_snia\_vs\_nonia} is based on the SuperNNova deep learning 
classifier \citep{Moller2019} and outputs a general $P(\mathrm{SNIa})$ 
score. \texttt{rf\_snia\_vs\_nonia} is a Random Forest classifier 
documented as targeting \textit{rising} SN Ia candidates; its scores 
are designed for early-phase, pre-maximum alerts rather than 
general-purpose SN Ia posteriors. We include both scores in our 
analysis but explicitly flag this distinction when interpreting the 
RF results in Sections~\ref{sec:results:fink} and~\ref{sec:discussion:fink}.

For each object we retrieved the most recent alert-level score available 
in the Fink database. We distinguish three outcomes: object found with 
scores (collected), object absent from Fink (not in ingestion window, 
treated as missing data), and API error (excluded). Absent objects are 
not coerced to zero — they are excluded from the Fink analysis entirely.

Of 1,606 objects queried, 1,368 (85.2\%) were present in the Fink 
database, with 237 objects absent due to falling outside Fink's alert 
ingestion window. Both classifiers returned scores for all 1,368 matched 
objects, yielding a comparative subsample of 644 SNIa, 480 SNII, 
166 SNIbc, 50 SLSN, and 28 TDE.

\subsection{Ground Truth and Selection Effects}
\label{sec:data:groundtruth}

We emphasise that our ground truth labels derive from independent 
spectroscopic observations, not from the broker classifications being 
evaluated. This avoids the circular validation present in studies that 
use broker-assigned labels to assess broker performance.

The BTS magnitude limit ($m < 18.5$) introduces a selection bias toward 
brighter, lower-redshift transients. Our sample therefore represents 
the regime in which spectroscopic follow-up is most feasible, which is 
precisely the regime most relevant to the deferral decisions this work 
motivates. We do not claim our calibration results generalise to 
fainter or higher-redshift populations. As the BTS selects transients 
for spectroscopic follow-up partly on the basis of brightness and 
observability, our sample likely over-represents well-behaved, 
unambiguous examples of each class. Calibration performance on the 
fainter, more ambiguous objects that will dominate LSST alert streams 
may differ; our results should be interpreted as an upper bound on 
broker performance in the bright regime.